\documentclass[10pt]{article}
\providecommand{\FIGDIR}{fig}
\newcommand{\DIGESTDATE}{26 August 2026}
% ===================== PM Research Watch - shared preamble =====================
\usepackage[T1]{fontenc}
\usepackage[utf8]{inputenc}
\usepackage{amsmath,amssymb}
\usepackage{tgpagella}
\usepackage{tgheros}
\renewcommand{\sfdefault}{qhv}
\renewcommand{\ttdefault}{lmtt}
\usepackage{microtype}
\usepackage[english]{babel}

\usepackage[a4paper,top=17mm,bottom=17mm,left=17mm,right=17mm,
            headsep=6mm,footskip=10mm]{geometry}
\usepackage{graphicx}
\usepackage{xcolor}
\usepackage{colortbl}
\usepackage{tikz}
\usetikzlibrary{calc}
\usepackage[most]{tcolorbox}
\usepackage{titlesec}
\usepackage{fancyhdr}
\usepackage{booktabs}
\usepackage{array}
\usepackage{longtable}
\usepackage{enumitem}
\usepackage{multicol}
\usepackage{ragged2e}
\usepackage[hidelinks]{hyperref}

% ------------------------------------------------------------------ palette
\definecolor{Ink}    {HTML}{14262E}
\definecolor{Deep}   {HTML}{0F4C5C}
\definecolor{Teal}   {HTML}{1F7A8C}
\definecolor{Sky}    {HTML}{4FA3C4}
\definecolor{Sage}   {HTML}{6FA287}
\definecolor{Amber}  {HTML}{E3A008}
\definecolor{Clay}   {HTML}{B5643C}
\definecolor{Coral}  {HTML}{D1495B}
\definecolor{Violet} {HTML}{6A5B9E}
\definecolor{Slate}  {HTML}{4A6572}
\definecolor{Paper}  {HTML}{FBF9F5}
\definecolor{Mist}   {HTML}{EEEAE2}
\definecolor{Rule}   {HTML}{D9D3C9}

\pagecolor{Paper}
\color{Ink}

\hypersetup{colorlinks=true, linkcolor=Deep, urlcolor=Teal, citecolor=Teal,
            pdftitle={PM Research Watch}, pdfauthor={Automated daily digest}}

% ------------------------------------------------------------------ headings
\titleformat{\section}
  {\sffamily\bfseries\large\color{Deep}}{}{0pt}
  {}[{\vspace{1.2mm}\color{Amber}\titlerule[1.1pt]}]
\titlespacing*{\section}{0pt}{5.5mm}{2.8mm}

\titleformat{\subsection}
  {\sffamily\bfseries\normalsize\color{Teal}}{}{0pt}{}
\titlespacing*{\subsection}{0pt}{3.6mm}{1.4mm}

\setlist[itemize]{leftmargin=4.2mm, itemsep=1.1mm, topsep=1.3mm, parsep=0pt,
                  label=\textcolor{Teal}{\footnotesize$\blacktriangleright$}}

% ------------------------------------------------------------------ header/footer
\pagestyle{fancy}
\fancyhf{}
\renewcommand{\headrulewidth}{0pt}
\renewcommand{\footrulewidth}{0pt}
\fancyhead[L]{\sffamily\scriptsize\color{Slate}PM RESEARCH WATCH\;\textcolor{Rule}{\textbar}\;\DIGESTDATE}
\fancyhead[R]{\sffamily\scriptsize\color{Slate}Particulate matter monitoring \& health effects}
\fancyfoot[C]{\sffamily\scriptsize\color{Slate}\thepage}

% ------------------------------------------------------------------ components
% metric tile
\newcommand{\metric}[3]{%
  \begin{tcolorbox}[enhanced, colback=white, colframe=Rule, boxrule=0.5pt,
    arc=1.6mm, left=2.2mm, right=2.2mm, top=1.8mm, bottom=1.8mm,
    borderline west={1.7mm}{0pt}{#3}, nobeforeafter, valign=center,
    height=17mm, width=\linewidth]
    \RaggedRight
    {\sffamily\bfseries\fontsize{17}{17}\selectfont\color{#3}#1}\\[0.5mm]
    {\sffamily\fontsize{6.5}{7.4}\selectfont\color{Slate}#2}
  \end{tcolorbox}}

% pull-quote / key insight box
\newtcolorbox{keybox}[1][]{enhanced, colback=Mist, colframe=Mist,
  arc=1.6mm, left=3.2mm, right=3.2mm, top=2.4mm, bottom=2.4mm,
  borderline west={1.7mm}{0pt}{Amber}, #1}

% subtopic band
\newcommand{\band}[2]{%
  \vspace{3.2mm}\noindent
  \begin{tcolorbox}[enhanced, colback=#2, colframe=#2, arc=1.2mm,
    left=2.6mm, right=2.6mm, top=1.1mm, bottom=1.1mm, boxrule=0pt,
    width=\linewidth, nobeforeafter]
    {\sffamily\bfseries\small\color{white}#1}
  \end{tcolorbox}\vspace{1.6mm}}

% one paper entry: {short cite}{title}{journal}{doi}{body}
\newcommand{\paperentry}[5]{%
  \par\vspace{1.9mm}\noindent
  {\sffamily\bfseries\footnotesize\color{Ink}#1}\;%
  {\sffamily\fontsize{7}{8}\selectfont\color{Slate}#3}\\[0.4mm]
  {\footnotesize\itshape\color{Deep}#2}\\[0.6mm]
  {\footnotesize #4}\\[0.3mm]
  {\fontsize{6.6}{7.6}\selectfont\color{Slate}\texttt{doi:}\,\href{https://doi.org/#5}{\texttt{\detokenize{#5}}}}
  \par}

% caption style
\newcommand{\figcap}[1]{\par\vspace{1.0mm}
  {\fontsize{7.2}{8.6}\selectfont\color{Slate}\textsf{\textbf{Fig.}}\; #1\par}\vspace{1.5mm}}

\providecommand{\permil}{\textperthousand}

\setlength{\parindent}{0pt}
\setlength{\parskip}{0pt}
\renewcommand{\arraystretch}{1.18}


\begin{document}

% ==================================================================== MASTHEAD
\thispagestyle{empty}
\begin{tikzpicture}[remember picture, overlay]
  \fill[Deep] (current page.north west) rectangle
      ($(current page.north east)+(0,-67mm)$);
  \fill[Amber] ($(current page.north west)+(0,-67mm)$) rectangle
      ($(current page.north east)+(0,-68.4mm)$);
  \foreach \x/\y/\r in {22/12/0.9, 41/26/1.5, 63/9/0.7, 88/31/1.1, 112/17/1.9,
                        138/29/0.8, 159/13/1.3, 178/54/1.0, 196/21/1.6,
                        30/44/0.6, 74/57/1.2, 125/48/0.9, 168/6/0.8, 191/59/0.7,
                        52/18/1.0, 100/60/0.6, 148/20/1.1, 66/49/1.4, 205/45/0.8}
    \fill[Sky, opacity=0.30] ($(current page.north west)+(\x mm,-\y mm)$) circle (\r mm);
\end{tikzpicture}

\vspace*{4mm}
{\sffamily\fontsize{7.6}{9}\selectfont\color{Sky}%
 AUTOMATED DAILY LITERATURE DIGEST\;\textbullet\;PhD RESEARCH WATCH\par}
\vspace{2.2mm}
{\sffamily\bfseries\fontsize{27}{29}\selectfont\color{white}%
 Particulate Matter\par}
{\sffamily\fontsize{27}{29}\selectfont\color{Sky}%
 Monitoring \& Health Effects\par}
\vspace{5.0mm}
{\sffamily\fontsize{8.4}{10}\selectfont\color{white}%
 \mbox{Issue 2026-08-26}\;\;\textcolor{Amber}{\textbar}\;\;\mbox{Entry window 26 Aug 2026}\;\;
 \textcolor{Amber}{\textbar}\;\;\mbox{14 in scope + 6 backfill, 36 out}\par}
\vspace{18mm}

% -------------------------------------------------------------- metric strip
\noindent
\begin{minipage}[t]{0.190\linewidth}\metric{14}{IN-WINDOW RECORDS\\(36 SCREENED OUT)}{Deep}\end{minipage}\hfill
\begin{minipage}[t]{0.190\linewidth}\metric{7}{CLUSTERS, EXACTLY\\2 RECORDS EACH}{Teal}\end{minipage}\hfill
\begin{minipage}[t]{0.190\linewidth}\metric{9}{INTERVAL-BEARING\\ESTIMATES}{Sky}\end{minipage}\hfill
\begin{minipage}[t]{0.190\linewidth}\metric{6}{OUT-OF-WINDOW\\CALIBRATION BACKFILL}{Amber}\end{minipage}\hfill
\begin{minipage}[t]{0.190\linewidth}\metric{2}{ABSTRACT-FREE\\DEPOSITS (14\%)}{Clay}\end{minipage}

\vspace{6mm}

% -------------------------------------------------------------- signal of day
\section*{\sffamily\bfseries\color{Deep}Signal of the day}
\vspace{-1.5mm}{\color{Amber}\rule{\linewidth}{1.1pt}}\vspace{2.5mm}

\begin{keybox}
{\sffamily\bfseries\small\color{Deep}The reference method that every automated PM monitor
is traceable to is itself humidity-sensitive --- and the EU's post-2030 uncertainty budget
does not leave room for that.}\\[1.4mm]
\footnotesize
An ISPRA-organised national interlaboratory comparison at Castel Romano put the
EN\,12341 gravimetric reference method for PM$_{10}$ and PM$_{2.5}$ in front of
\textbf{18 Italian regional environmental agencies} simultaneously. Systematic deviations
between laboratories were driven \textbf{primarily by relative humidity and
precipitation}, which compromise PM thermodynamic and mass stability during handling and
conditioning, with a smaller filter-matrix contribution. Gravimetry is the sole
traceability source for every automated measurement system, so this bias propagates
downward into the entire monitoring chain --- and \textbf{Directive (EU) 2024/2881
sharply tightens AMS residual uncertainty budgets for daily and annual PM$_{10}$ and
PM$_{2.5}$ from 2030}, at concentrations that are themselves declining. The uncomfortable
implication for low-cost sensor work: the humidity artefact that every LCS calibration
paper treats as a sensor problem is partly present in the reference the sensor is being
calibrated \emph{against}. Co-location residuals attributed to sensor RH response may
contain a reference-side component that no amount of sensor-side modelling can remove.
\end{keybox}

\vspace{3mm}

% -------------------------------------------------------------- exec brief
\section*{\sffamily\bfseries\color{Deep}What changed today}
\vspace{-1.5mm}{\color{Amber}\rule{\linewidth}{1.1pt}}\vspace{2.5mm}

\footnotesize
\begin{itemize}\itemsep2pt
\item \textbf{The cleanest component-resolved result this watch has carried.}
Mao et al.\ ran 60,789 consecutive gastroscopies against gridded daily PM$_{2.5}$ and its
five chemical components. Total PM$_{2.5}$ gave \textbf{IRR 1.076 (95\,\% CI
1.024--1.132)} per IQR for endoscopically confirmed erosive gastric lesions, but the
signal sits entirely in the \textbf{secondary inorganic fraction --- ammonium 1.082
(1.026--1.141) and sulfate 1.070 (1.024--1.118)} --- while nitrate, organic matter and
black carbon are all flat at 1.012 with intervals spanning the null. Six estimates, one
study, one exposure model, one outcome ascertainment: the internal comparison is as
clean as this literature gets, and it says mass is the wrong metric.

\item \textbf{Ultrafine particles get a mediation pathway.} Wang et al.\ modelled UFP at
high spatiotemporal resolution for 7,569 Shanghai pregnancies and found per-IQR relative
risks of 1.15--1.23 for hypertensive disorders of pregnancy and 1.29--1.41 for
pre-eclampsia. First-trimester homocysteine mediated 2.51--8.87\,\% of the UFP--pre-eclampsia
association (\textbf{Hcy--PE RR 1.49, 1.04--2.13}), and the association held only in the
low-vitamin-B$_{12}$ stratum. Mediated fractions under 10\,\% are easy to over-read;
the B$_{12}$ modification is the more actionable finding.

\item \textbf{A large-cohort PM/BPH result with a stated threshold.} Wei et al.\ report
2.18\,\% higher BPH risk per 1\,$\mu$g/m$^3$ PM$_{2.5}$ across 123,037 UK Biobank men,
replicated in a western-China cohort at 2.77\,\%, with a restricted-cubic-spline
\textbf{risk threshold at 9.83\,$\mu$g/m$^3$} --- essentially the WHO interim target and
the same $\ge$9\,$\mu$g/m$^3$ breakpoint the AF literature landed on last month. Thirty-four
lipid metabolites are proposed as mediators, which at this cohort size is exploratory.

\item \textbf{A dietary effect modification worth scepticism.} Kessler et al.\ report
that a first-trimester IQR increase in PM was associated with a 0.08-unit
\emph{decrease} in birth-weight z-score at mean $\alpha$-linolenic-acid intake but a
\textbf{1.12-unit increase (0.19--2.05) at twice mean ALA intake} ($p_{\mathrm{int}}=0.04$).
A sign reversal of that magnitude on $n=297$, driven by the upper tail of a
food-frequency-questionnaire distribution, is far more likely to be sparse-data bias than
a real protective interaction. Reported here because it will be cited.

\item \textbf{Six-record coverage backfill on low-cost sensor calibration} (see the
sidebar). A Consensus topical sweep surfaced six 2026 calibration papers deposited
between January and June --- before the Crossref by-ISSN leg existed --- that no leg of
this pipeline has ever seen. They are \emph{not} counted in the metric strip, the figures
or the register, which describe the 26 August window only.

\item \textbf{Two abstract-free deposits, both squarely on topic.} Condensable
particulate matter in subzero remote Canadian mountain communities
(\textit{Environ Int}) and source-resolved PM$_{2.5}$ oxidative potential in a
shipbuilding complex (\textit{Atmos Pollut Res}) arrived with title and DOI only. The
Elsevier abstract backlog now spans four weeks and both go on the retry list at the top.
\end{itemize}

\clearpage

% ==================================================================== ANALYTICS
\section*{\sffamily\bfseries\color{Deep}Corpus at a glance}
\vspace{-1.5mm}{\color{Amber}\rule{\linewidth}{1.1pt}}\vspace{3mm}

\noindent\includegraphics[width=\linewidth]{\FIGDIR/f1_subtopics.png}
\figcap{Records per subtopic cluster. \textbf{Seven clusters, exactly two records each} ---
a flat distribution this watch has not produced before, and an artefact of a single
\textit{Environment \& Health} issue release landing four unrelated endpoint papers in
one day rather than of any convergence in the field. \textit{Cardiovascular \&
metabolic}, \textit{Neuro / mental health} and \textit{Occupational \& indoor} are empty.}

\vspace{2mm}
\noindent\includegraphics[width=\linewidth]{\FIGDIR/f2_design_endpoint.png}
\figcap{Left --- study architecture as binned by the coarse grouper.
\textit{Measurement campaign} 4, \textit{Observational --- cohort} 3,
\textit{Observational --- acute} 2, \textit{Experimental / toxicology} 2, and one each for
\textit{Review / synthesis}, \textit{Ecological} and \textit{Modelling / inventory}.
The measurement-campaign bar is heterogeneous: it holds an interlaboratory gravimetric
comparison, an electron-microscopy characterisation and two abstract-free field records
whose design is inferred from the title. Right --- health endpoint. \textbf{Five records
carry no human health endpoint} (two instrument, two exposure, one policy); the
remaining nine endpoints are all singletons.}

\vspace{2mm}
\noindent\includegraphics[width=\linewidth]{\FIGDIR/f3_metric_geography.png}
\figcap{Left --- particle metric. \textit{PM$_{2.5}$ only} 7 and
\textit{PM$_{2.5}$ + PM$_{10}$ jointly} 5; \textbf{one record measures ultrafine particle
number} (Wang et al., the Shanghai pregnancy cohort) and one is unspeciated. Note that
Mao et al.\ is binned as PM$_{2.5}$-only although its actual exposure is five resolved
chemical components --- the grouper has no speciation bin, which is a real limitation
given that speciation is where this window's signal lives. Right --- geographic
provenance: China 4, North America 3, multi-region syntheses 3, Europe 2, and one each
from the Middle East \& North Africa (Saudi Arabia) and East Asia outside China
(South Korea). Nothing from Africa south of the Sahara, South Asia or Latin America.}

\vspace{2mm}
\noindent\includegraphics[width=\linewidth]{\FIGDIR/f4_heatmap.png}
\figcap{Subtopic $\times$ study-architecture cross-tabulation. With 14 records spread over
seven clusters and seven architectures the matrix is almost entirely singletons. Four
cells hold 2 --- \textit{Sensing} and \textit{Exposure assessment} both $\times$
\textit{Measurement campaign}, \textit{Reproductive \& developmental} $\times$
\textit{Observational --- cohort}, and \textit{Mechanistic toxicology} $\times$
\textit{Experimental / toxicology} --- and every other occupied cell holds 1. The two
architecture columns that carry no health endpoint anywhere
(\textit{Measurement campaign}, \textit{Modelling / inventory}) account for 5 of the 14
records. Read as a sparsity map rather than as structure.}

\vspace{2mm}
\noindent\includegraphics[width=\linewidth]{\FIGDIR/f5_forest.png}
\figcap{Nine interval-bearing estimates from three studies. \textbf{The six Mao et al.\
estimates are mutually comparable} --- same cohort, same exposure model, same outcome,
same per-IQR contrast --- and are the only set in this watch's history for which a
component-versus-component reading is legitimate. The two COPD odds ratios are on a
per-unit rather than per-IQR contrast and are not comparable to them; note also that the
publisher's indexing stripped the subscripts from the source abstract, so the assignment
of finer and coarser PM fraction follows the order in which the abstract lists them and
should be confirmed against the full text. The homocysteine relative risk is a
\emph{mediator--outcome} estimate, not an exposure--outcome estimate, and belongs on this
axis only in the weak sense that it shares a ratio scale.}

\vspace{2mm}
\noindent\includegraphics[width=\linewidth]{\FIGDIR/f6_lifecourse.png}
\figcap{Life-course windows addressed, counted only where a record's study population is
explicitly anchored to a stage. Six records qualify. \textbf{The preconception/in-utero
window is occupied for the first time since 22 August} (3 records, after three
consecutive empty issues), and the 65+ window carries the Kashi COPD series, whose
population is largely older adults. Paediatric and adolescent remain empty --- the
paediatric window has now been empty for four consecutive issues.}

\clearpage

% ==================================================================== BACKFILL
\section*{\sffamily\bfseries\color{Deep}Coverage backfill: low-cost sensor calibration}
\vspace{-1.5mm}{\color{Amber}\rule{\linewidth}{1.1pt}}\vspace{2.5mm}

\begin{keybox}
{\sffamily\bfseries\footnotesize\color{Deep}Six 2026 papers this pipeline never saw ---
out of window, not counted anywhere else in this issue}\\[1.2mm]
{\fontsize{7.2}{8.8}\selectfont
A Consensus topical sweep on low-cost sensor calibration, drift, co-location and humidity
correction returned six 2026 records, \textbf{all six absent from
\texttt{seen.json}}. Crossref \texttt{created} dates run 23 January to 21 June 2026 ---
that is, before the by-ISSN sensing leg was built on 22 August, and in journals
(\textit{Meas Sci Technol}, \textit{Aerosol Sci Technol}) that leg still does not track.
This is a genuine historical coverage hole, not a duplicate. They are recorded in
\texttt{seen.json} so they will not re-ship, and both untracked ISSNs have been added to
the harvester.\par}
\end{keybox}

\vspace{2mm}
\footnotesize
\begin{itemize}\itemsep2pt
\item \textbf{Ginsburg et al., \textit{Aerosol Sci Technol}} ---
\textbf{1,013 PurpleAir sensors} co-located within 500\,m of reference monitors on six
continents, January 2023 to March 2025. Raw PurpleAir overestimates PM$_{2.5}$ in most
regions and fails EPA sensor performance targets; \textbf{US-derived corrections do not
transfer outside North America}; a global quadratic helps in Australia/New Zealand and
hurts elsewhere; \textbf{regional calibrations meet EPA 24\,h thresholds in multiple
regions and lift 24\,h AQI category correctness to $\ge$77\,\%}. The single most
consequential record in this backfill for anyone running a non-US network.
\texttt{10.1080/02786826.2026.2676293}

\item \textbf{Kaur et al., \textit{Atmos Meas Tech}} --- \textbf{109 PMS5003/6003 sensors}
at or near 75 US EPA sites, January 2017 to May 2025. Two thresholds built from relative
humidity and the sensor's own coarse-to-submicron mass ratio detect dust-dominated
periods \emph{from sensor data alone}, using the trailing 336 hourly values that
PurpleAir already archives. Applied as a non-linear correction to \texttt{pm2.5\_alt},
this \textbf{cut mean bias error by ${\sim}50\,\%$ on 97 sensors and RMSE by
${\sim}30\,\%$ on 84}. Directly implementable with no additional hardware.
\texttt{10.5194/amt-19-1077-2026}

\item \textbf{Gupta et al., \textit{Aerosol Air Qual Res}} --- PA-II against FEM in
Raleigh and Delhi, 2--4 week deployments chosen to span contrasting aerosol loadings.
Calibrated output reaches \textbf{within 12\,\% mean absolute bias hourly and 6\,\%
daily}. The authors conclude pre-deployment regional calibration is a precondition for
scientific use, which is the same conclusion Ginsburg et al.\ reach from a global
dataset. \texttt{10.1007/s44408-026-00146-x}

\item \textbf{Zhivkov et al., \textit{Atmosphere}} --- the calibration-transfer problem
stated properly: what happens to a machine-learning correction applied to a sensor that
is nowhere near a reference station. 24 months, Sofia, 5 reference stations and 22
sensors of which 4 are co-located. Random Forest reached $R^2\in(0.53,0.75)$, and
leave-one-station-out cross-validation yielded \textbf{uncertainty growth of
3.84--5.62\,\% per km} from the nearest reference (significant for PM$_{10}$ and O$_3$).
\textbf{Co-location sites 6\,km apart hold midpoint uncertainty growth under 30\,\%},
which is an actionable network-siting rule. \texttt{10.3390/atmos17040335}

\item \textbf{Qian et al., \textit{Atmos Meas Tech}} --- AutoML calibration for indoor
PM$_{2.5}$, staged through an intermediate drift-correction reference and fitted
separately for clean-air and pollution-event concentration ranges. Chamber evaluation on
two sensor models gave \textbf{$R^2>0.90$ with NRMSE and sMAPE roughly halved}. The
split-range design is the transferable part; chamber validation limits the field claim.
\texttt{10.5194/amt-19-603-2026}

\item \textbf{Feng et al., \textit{Meas Sci Technol}} --- review of machine-learning,
physics-based and hybrid physics-informed calibration. Conclusion worth quoting:
\textbf{no single strategy is universally optimal} --- ML wins on predictive performance,
physics on interpretability and physical consistency, hybrids on transferability --- and
non-linear models reaching sub-10\,$\mu$g\,m$^{-3}$ error still face embedded-deployment,
training-data-representativeness and region-specificity barriers.
\texttt{10.1088/1361-6501/ae6298}
\end{itemize}

\clearpage

% ==================================================================== DIGEST
\section*{\sffamily\bfseries\color{Deep}Paper-level digest}
\vspace{-1.5mm}{\color{Amber}\rule{\linewidth}{1.1pt}}\vspace{1mm}

{\fontsize{7.4}{9}\selectfont\color{Slate}Ordered by cluster; the 26 August window only.
The six backfill records are covered in the preceding section and are not repeated here.\par}

\band{Sensing, forecasting \& instrumentation\hfill 2 records}{Amber}

\paperentry{Anon et al.\ 2026}
{Metrological consistency of the EN 12341 gravimetric reference method}
{Atmosphere}
{National interlaboratory comparison run by ISPRA at Castel Romano with \textbf{18
Italian regional environmental agencies} sampling side by side, motivated by the absence
of Certified Reference Materials for PM and by known sampler-design discrepancies.
Systematic between-laboratory deviations were driven mainly by \textbf{relative humidity
and precipitation} acting on PM thermodynamic and mass stability, with a minor
filter-matrix effect. Because gravimetry is the only traceability source for automated
measurement systems, the bias propagates into every AMS calibrated against it --- and
Directive (EU) 2024/2881 tightens AMS residual uncertainty budgets sharply from 2030,
at concentrations that are falling. No numeric bias magnitudes appear in the abstract,
which is the main gap; the full text is needed to size the effect. Crossref lists no
first author.}
{10.3390/atmos17090828}

\paperentry{Alharbi et al.\ 2026}
{Morphology, surface chemistry and mineralogy of coastal PM$_{2.5}$, Jeddah}
{Atmosphere}
{Three coastal sites over summer 2017 characterised by SEM, EDS mapping, XPS and XRD,
linking bulk concentration (\textbf{22.2, 18.9 and 14.2\,$\mu$g/m$^3$} north to south)
to particle morphology, surface oxidation state and crystalline phase. The authors are
unusually candid about their own artefact: \textbf{borosilicate glass-fibre substrates
dominate the SEM field}, so aerosol is resolved only as discrete sub-micrometre particles
and agglomerates decorating the fibres, and the morphological interpretation is framed
accordingly. That admission is worth more than the result --- filter-substrate choice
constrains what electron microscopy of a PM sample can establish, and most papers in this
genre do not say so.}
{10.3390/atmos17090830}

\band{Exposure assessment \& modelling\hfill 2 records}{Deep}

\paperentry{Anon et al.\ 2026b}
{Condensable particulate matter and severe winter PM$_{2.5}$ in remote Canadian mountain communities}
{Environ Int}
{\textbf{Metadata-only deposit} --- title and DOI without an abstract. The topic is
directly relevant and under-covered: condensable PM forms downstream of the stack as
vapour cools, is invisible to filter-based source testing conducted at temperature, and
in subzero valley-inversion conditions can dominate ambient PM$_{2.5}$ in communities
with no industrial source of any size. If the finding is what the title implies, it also
bears on optical sensing --- semi-volatile condensable material has different scattering
and hygroscopic behaviour from the combustion soot most correction factors assume.
Carried at tier C and flagged as the top abstract retry.}
{10.1016/j.envint.2026.110472}

\paperentry{Anon et al.\ 2026c}
{Source-specific PM$_{2.5}$ oxidative potential in a national shipbuilding complex}
{Atmos Pollut Res}
{\textbf{Metadata-only deposit.} The title asserts that metal processing and combustion
dominate source-apportioned oxidative potential, which if supported is another instance
of the recurring result that OP-weighted and mass-weighted source contributions rank
differently. No apportionment method, OP assay, site count or period is available.
Tier C, second on the abstract retry list.}
{10.1016/j.apr.2026.103189}

\clearpage

\band{Reproductive \& developmental\hfill 2 records}{Sky}

\paperentry{Wang et al.\ 2026}
{Ultrafine particles, hypertensive disorders of pregnancy, and homocysteine mediation}
{Environ Health (Wash)}
{Shanghai birth cohort, \textbf{7,569 pregnancies}, 2017--2018, with individual UFP
exposure from a high-spatiotemporal-resolution prediction model at residential address
and serum homocysteine, folate and B$_{12}$ measured in early pregnancy. Poisson
regression with robust error variance gave per-IQR \textbf{RR 1.15--1.23 for
hypertensive disorders of pregnancy and 1.29--1.41 for pre-eclampsia} across the
prepregnancy, first- and second-trimester windows. Homocysteine was associated with
pre-eclampsia (\textbf{RR 1.49, 1.04--2.13}) and mediated 8.87\,\% and 2.51\,\% of the
prepregnancy and first-trimester associations. Stratification showed the UFP--PE
association only in the low-B$_{12}$ group. Two cautions: the exposure is modelled, not
measured, so UFP--PM$_{2.5}$ collinearity is not resolvable here; and mediated fractions
this small are within the range that measurement error in the mediator can manufacture.}
{10.1021/envhealth.5c00760}

\paperentry{Kessler et al.\ 2026}
{Prenatal traffic-related air pollution, dietary PUFA intake and birth weight}
{Environ Health (Wash)}
{Atlanta Maternal--Child Cohort, \textbf{297 pregnant African American women},
first-trimester traffic-related CO, PM$_{2.5}$ and NO$_x$ crossed with dietary
polyunsaturated fatty acid intake in a joint mixture model. At mean $\alpha$-linolenic
acid intake, a per-IQR PM increase was associated with a \textbf{0.08-unit decrease
(0.02--0.18) in birth-weight-for-gestational-age z-score}; at twice mean ALA intake the
same increase was associated with a \textbf{1.12-unit increase (0.19--2.05)},
$p_{\mathrm{int}}=0.04$, with an almost identical pattern for linoleic acid. The authors
read this as PUFA mitigation. A cleaner reading is that a 14-fold sign-reversing
interaction estimated on 297 women in the sparse upper tail of a food-frequency
distribution is the signature of sparse-data bias. Worth tracking because the
mechanistic story is plausible and the result will be cited widely.}
{10.1021/envhealth.5c00490}

\band{Other clinical endpoints\hfill 2 records}{Slate}

\paperentry{Mao et al.\ 2026}
{Component-specific PM$_{2.5}$ and endoscopically confirmed erosive gastric lesions}
{Front Public Health}
{Consecutive gastroscopies at Yixing People's Hospital, 1 Jan 2024 to 31 Dec 2025:
\textbf{60,789 individuals, 14,325 with erosive gastric mucosal lesions}, against gridded
daily PM$_{2.5}$ and five chemical components in generalized additive models with
attributable-burden analysis. Per IQR, \textbf{total PM$_{2.5}$ IRR 1.076
(1.024--1.132)}; after FDR correction only \textbf{ammonium 1.082 (1.026--1.141)} and
\textbf{sulfate 1.070 (1.024--1.118)} survived, with \textbf{nitrate, organic matter and
black carbon all at 1.012} and intervals crossing the null. Endoscopic ascertainment
removes the symptom-reporting bias that weakens most digestive air-pollution studies.
The limitation is real and stated: the five components are strongly collinear and
share one exposure model, so their independent contributions are not identified. Still
the sharpest secondary-inorganic-versus-carbonaceous contrast available.}
{10.3389/fpubh.2026.1891236}

\paperentry{Wei et al.\ 2026}
{Air pollution, lipid mediators and benign prostatic hyperplasia}
{Environ Health (Wash)}
{UK Biobank, \textbf{123,037 men}, with a western-China cohort replication.
\textbf{2.18\,\% higher BPH risk per 1\,$\mu$g/m$^3$ PM$_{2.5}$} ($p=0.0029$) in UK
Biobank and \textbf{2.77\,\%} in the replication; mixture exposure by WQS and qgcomp gave
3.53\,\% and 2.20\,\% per unit for BPH and 6.78\,\% and 9.17\,\% for lower urinary tract
symptoms. Restricted cubic splines put the \textbf{BPH risk threshold at
9.83\,$\mu$g/m$^3$}, and joint high-PM/high-fat-diet exposure carried 16\,\% higher risk
than low/low. Thirty-four lipid metabolites are nominated as mediators. Two problems:
BPH ascertainment in UK Biobank is by linked diagnosis and is sensitive to
healthcare-contact patterns that correlate with urbanicity and therefore with PM$_{2.5}$;
and 34 mediators from an untargeted panel without a stated multiplicity correction is
hypothesis generation, not mediation analysis.}
{10.1021/envhealth.5c00806}

\clearpage

\band{Respiratory \& allergic\hfill 2 records}{Sage}

\paperentry{Li et al.\ 2026}
{Air pollution and COPD hospitalisation in Kashi, Xinjiang}
{Int J Chron Obstruct Pulmon Dis}
{Time-stratified case-crossover on \textbf{68,482 COPD patients} in Kashi, 2012--2022,
with daily pollutant concentrations matched to residential address rather than to a
single city monitor. Strongest PM association at \textbf{lag day 4}: odds ratios
\textbf{1.015 (1.002--1.032)} and \textbf{1.008 (1.002--1.016)} for the two PM fractions,
with NO$_2$ and CO both at 1.034. Young patients, those with prior severe exacerbations,
farmers, smokers and cool-season admissions were the susceptible strata. The honest part
is the sensitivity analysis: \textbf{excluding COVID-19-period hospitalisations left the
PM and NO$_2$ associations barely significant and CO non-significant}, which means a
material share of the headline estimate rides on pandemic-era admission patterns.
The publisher's indexing stripped every subscript from the abstract, so which PM fraction
is which must be read from listing order. Kashi sits in a severe dust-storm corridor,
making this one of the few large COPD series in a mineral-dust-dominated aerosol.}
{10.2147/COPD.S617521}

\paperentry{Wang \& Li 2026}
{Air pollution and atopic dermatitis}
{Front Public Health}
{Narrative review of PubMed, Web of Science and CNKI, April 2021 to December 2025,
restricted to peer-reviewed work with impact factor $\ge$3.0 --- a screening criterion
that is a proxy for nothing and silently excludes most dermatology and most regional
epidemiology. Covers epidemiology, skin-barrier and microbiome mechanisms, climatic and
geographical disparity, and prevention. No quantitative synthesis and no effect
estimates. Carried at tier C as a citation anchor for the cutaneous endpoint, which is
otherwise absent from this watch.}
{10.3389/fpubh.2026.1854212}

\band{Mechanistic toxicology\hfill 2 records}{Violet}

\paperentry{Liu et al.\ 2026}
{PM$_{2.5}$-induced pulmonary vascular remodelling via MSC-mediated endothelial senescence}
{Environ Health (Wash)}
{Real-ambient PM$_{2.5}$ exposure in mice plus a pulmonary arterial endothelial cell
model. Exposure raised pulmonary artery pressure and arterial wall thickness in vivo and
induced senescence and mitochondrial dysfunction in mouse pulmonary arterial endothelial
cells. Bioinformatic screening nominated the transcription factor MSC as the regulator;
\textbf{MSC overexpression rescued mitochondrial-dysfunction-associated senescence}, and
MSC was shown to act through the downstream target NDUFB3, a complex I subunit ---
which makes the mitochondrial phenotype mechanistically coherent rather than merely
correlative. \textbf{Real-ambient rather than instilled exposure} keeps the dosimetry
defensible. No concentration, duration or particle characterisation appears in the
abstract, so comparability to other real-ambient systems cannot be judged.}
{10.1021/envhealth.5c00595}

\paperentry{TaRGET II Consortium 2026}
{Lifecourse sex-specific molecular response to early-life toxicant exposure}
{bioRxiv (preprint)}
{Consortium-scale design: mice exposed from two weeks preconception through gestation and
lactation to arsenic, lead, tributyltin, BPA, DEHP, TCDD or \textbf{PM$_{2.5}$}, with
\textbf{1,025 omics datasets} generated from liver across three life stages. The
headline is a significant sex-biased transcriptomic and epigenomic response to a common
exposure. For this watch the value is comparative --- PM$_{2.5}$ is run alongside six
reference toxicants under one protocol, which is the only way to judge whether a PM
molecular signature is distinctive or generic. Liver-only, and mouse inhalation dosimetry
does not scale simply to humans. Unreviewed.}
{10.64898/2026.08.20.746014}

\band{Burden, policy \& mitigation\hfill 2 records}{Coral}

\paperentry{Anon et al.\ 2026d}
{Lifetime risk of respiratory-system cancers across 185 countries}
{Prev Med Rep}
{GLOBOCAN 2022 combined with UN life tables under an adjusted-for-multiple-primaries
framework. \textbf{Global lifetime risk of developing lung cancer 4.92\,\% and of dying
from it 4.08\,\%}, against 0.33\,\% and 0.20\,\% for laryngeal cancer, with a steep
gradient from \textbf{0.34\,\% in low-HDI to 5.45\,\% in very-high-HDI countries}. All
four pollutants were positively associated with lifetime risk in univariable linear
regression, PM$_{2.5}$ strongest. Two hard limits: the HDI gradient in lung cancer
incidence is dominated by historical smoking prevalence and by diagnostic capacity, and
country-level pollution from \textbf{OpenWeather} is a consumer weather API rather than a
validated exposure product. Univariable ecological regression across 185 countries
cannot separate PM from GDP, smoking or registry completeness. Tier B for the burden
estimates; the pollution association should not be cited as evidence of effect.}
{10.1016/j.pmedr.2026.103602}

\paperentry{Anon et al.\ 2026e}
{A Harmonization Index for industrial PM emission regulations}
{Res Sq (preprint)}
{A Regulatory Coding Framework converting legally binding documents into standardized
scores across \textbf{21 indicators in three dimensions} --- regulatory stringency,
institutional design and enforcement capacity --- aggregated into a composite
Harmonization Index, applied to the EU, US, China and Uzbekistan as exemplars of
technology-based, source-category-based, hybrid and dispersion-based regulatory
architectures. The problem is genuine: industrial PM limits are not comparable across
jurisdictions because they are not the same kind of object. The method is documentary
coding of what is written, which measures statutory text and not emissions; the
enforcement-capacity dimension is the one that would need external validation and is the
one hardest to score from documents. Unreviewed, four jurisdictions, tier C.}
{10.21203/rs.3.rs-10491363/v1}

\clearpage

% ==================================================================== TABLE
\section*{\sffamily\bfseries\color{Deep}Corpus register}
\vspace{-1.5mm}{\color{Amber}\rule{\linewidth}{1.1pt}}\vspace{2.5mm}

{\fontsize{6.7}{8.4}\selectfont
\setlength{\tabcolsep}{2.4pt}
\begin{longtable}{@{}>{\RaggedRight}p{25mm}>{\RaggedRight}p{41mm}>{\RaggedRight}p{27mm}%
>{\RaggedRight}p{23mm}>{\RaggedRight}p{18mm}>{\RaggedRight}p{15mm}@{}}
\toprule
\textsf{\textbf{First author}} & \textsf{\textbf{Focus}} & \textsf{\textbf{Journal}} &
\textsf{\textbf{Design}} & \textsf{\textbf{Metric}} & \textsf{\textbf{Region}}\\
\midrule
\endfirsthead
\toprule
\textsf{\textbf{First author}} & \textsf{\textbf{Focus}} & \textsf{\textbf{Journal}} &
\textsf{\textbf{Design}} & \textsf{\textbf{Metric}} & \textsf{\textbf{Region}}\\
\midrule
\endhead
\multicolumn{6}{@{}l@{}}{\cellcolor{Mist}\textsf{\textbf{\textcolor{Deep}{Neuro / mental health}}}}\\[0.6mm]
\href{https://doi.org/10.1001/jamanetworkopen.2026.25015}{Dales} & Wildfire-sourced PM$_{2.5}$ and ED visits for migraine and other primary headache & \textit{JAMA Netw Open} & Case-crossover & PM$_{2.5}$ (wildfire) & Canada\\
\href{https://doi.org/10.1093/ageing/afag215}{Glover} & Umbrella meta-review of urban environmental exposures and cognitive health & \textit{Age Ageing} & Umbrella review / meta-analysis & PM$_{2.5}$ & Global\\
\href{https://doi.org/10.1016/j.nsa.2026.107019}{Goossens} & In utero air pollution exposure and pre-/postnatal brain development & \textit{Neurosci Appl} & Narrative review & PM$_{2.5}$ / PM$_{10}$ / PM$_{\mathrm{coarse}}$ & Global\\
\href{https://doi.org/10.1016/j.isci.2026.116617}{Ma} & PM$_{2.5}$-induced neurotoxicity and depression: exposure risks and mechanisms & \textit{iScience} & Narrative review & PM$_{2.5}$ + components & Global\\
\href{https://doi.org/10.3760/cma.j.cn112150-20250804-00752}{Pan} & Maternal PM exposure and offspring neurodevelopment: effects and mechanisms & \textit{Zhonghua Yu Fang Yi Xue Za Zhi} & Narrative review & PM (bulk) & China\\
\addlinespace[0.9mm]
\multicolumn{6}{@{}l@{}}{\cellcolor{Mist}\textsf{\textbf{\textcolor{Teal}{Cardiovascular \& metabolic}}}}\\[0.6mm]
\href{https://doi.org/10.3760/cma.j.cn112338-20251120-00843}{Huang} & PM$_{2.5}$ and dyslipidaemia in rural Chinese children; dietary effect modification & \textit{Zhonghua Liu Xing Bing Xue Za Zhi} & Cross-sectional (multistage) & PM$_{2.5}$ & China\\
\href{https://doi.org/10.32394/pe/218966}{Kazmierska} & Air pollution and cardiovascular disease: review of current data & \textit{Przegl Epidemiol} & Narrative review & PM$_{2.5}$ / PM$_{10}$ & Global\\
\href{https://doi.org/10.1007/s00125-026-06814-2}{Xi} & Air pollution, built environment and temperature vs. prediabetes phenotypes & \textit{Diabetologia} & Prospective cohort & PM$_{2.5}$ & Germany\\
\href{https://doi.org/10.1161/JAHA.125.046917}{Yang} & Short-term PM$_{2.5}$ and hospitalisation/ED use in atrial fibrillation & \textit{J Am Heart Assoc} & Case-crossover & PM$_{2.5}$ (wildfire-influenced) & USA\\
\addlinespace[0.9mm]
\multicolumn{6}{@{}l@{}}{\cellcolor{Mist}\textsf{\textbf{\textcolor{Sky}{Reproductive \& developmental}}}}\\[0.6mm]
\href{https://doi.org/10.1016/j.jes.2025.10.043}{Chen} & Prenatal air pollution, fetal growth trajectory phenotypes and child neurodevelopment & \textit{J Environ Sci (China)} & Prospective cohort & PM$_{2.5}$ / PM$_{10}$ & China\\
\href{https://doi.org/10.1016/j.ecoenv.2026.120566}{Huang} & Joint prenatal bisphenol + air-pollutant mixture exposure and preschool IQ & \textit{Ecotoxicol Environ Saf} & Prospective cohort (mixtures) & PM$_{2.5}$ / PM$_{10}$ & China\\
\href{https://doi.org/10.1007/s40572-026-00552-8}{Peterson} & Ambient air pollution, PAHs and pubertal development: critical review & \textit{Curr Environ Health Rep} & Critical review & PM$_{2.5}$ components & Global\\
\href{https://doi.org/10.1016/j.envres.2026.125313}{Xie} & Long-term PM$_{2.5}$ and timing of pubertal development in boys and girls & \textit{Environ Res} & Prospective cohort & PM$_{2.5}$ & China\\
\addlinespace[0.9mm]
\multicolumn{6}{@{}l@{}}{\cellcolor{Mist}\textsf{\textbf{\textcolor{Sage}{Respiratory \& allergic}}}}\\[0.6mm]
\href{https://doi.org/10.1007/s00484-026-03289-z}{Bayar Muluk} & Climate change and otolaryngology: airway disease impacts and emission reduction & \textit{Int J Biometeorol} & Narrative review & PM$_{2.5}$ & Global\\
\href{https://doi.org/10.4103/mgr.MEDGASRES-D-25-00164}{Liao} & Short-term air pollution, respiratory symptoms and lung function in asthma & \textit{Med Gas Res} & Panel study & PM$_{2.5}$ / PM$_{10}$ & China\\
\href{https://doi.org/10.3390/epidemiologia7040097}{Sanchez de Prada} & Ambient PM and asthma-coded admissions in children 0-3 y, Spain & \textit{Epidemiologia} & Ecological panel & PM$_{2.5}$ / PM$_{10}$ & Spain\\
\href{https://doi.org/10.1016/j.jhazmat.2026.143076}{Yang} & Allergenic vs. toxic risk of airborne seed hairs compared with ambient PM$_{2.5}$ & \textit{J Hazard Mater} & Field measurement & PM$_{2.5}$ + bioaerosol & China\\
\addlinespace[0.9mm]
\multicolumn{6}{@{}l@{}}{\cellcolor{Mist}\textsf{\textbf{\textcolor{Violet}{Mechanistic toxicology}}}}\\[0.6mm]
\href{https://doi.org/10.1016/j.jes.2025.12.074}{Ding} & PM$_{2.5}$ promotes vascular calcification via miR-27b-5p/HES1 signalling & \textit{J Environ Sci (China)} & In vitro + organoid + murine & PM$_{2.5}$ & China\\
\href{https://doi.org/10.1007/s00360-026-01698-5}{Maraschi} & Osmoregulatory disruption and metal bioaccumulation from settleable atmospheric PM & \textit{J Comp Physiol B} & Controlled exposure (ecotox) & Settleable APM & Brazil\\
\href{https://doi.org/10.1016/j.jes.2025.11.018}{Wu} & Macrophage NLRP3 inflammasome in real-ambient PM-induced hepatic glucose dysregulation & \textit{J Environ Sci (China)} & Real-ambient murine + co-culture & Real-ambient PM & China\\
\addlinespace[0.9mm]
\multicolumn{6}{@{}l@{}}{\cellcolor{Mist}\textsf{\textbf{\textcolor{Clay}{Occupational \& indoor}}}}\\[0.6mm]
\href{https://doi.org/10.3389/fmedt.2026.1796828}{Chen} & Nested suction instrument to cut intraoperative surgical-smoke PM$_{2.5}$ & \textit{Front Med Technol} & Bench / device experiment & PM$_{2.5}$ (surgical smoke) & China\\
\href{https://doi.org/10.1093/annweh/wxag055}{Cole} & Determinants of respirable crystalline silica exposure in Australian tunnelling & \textit{Ann Work Expo Health} & Observational exposure study & Respirable dust / RCS & Australia\\
\href{https://doi.org/10.1016/j.jhazmat.2026.143081}{Ding} & Stage-resolved hair Zn-isotope dynamics as a firefighting combustion-exposure tracer & \textit{J Hazard Mater} & Repeated-measures biomarker & Combustion-derived PM & China\\
\href{https://doi.org/10.3390/jox16040119}{Ramsden} & Xenobiotic hazards in aircraft cabin air & \textit{J Xenobiot} & Review + risk economics & Ultrafine particles & Global\\
\addlinespace[0.9mm]
\multicolumn{6}{@{}l@{}}{\cellcolor{Mist}\textsf{\textbf{\textcolor{Amber}{Sensing, forecasting \& instrumentation}}}}\\[0.6mm]
\href{https://doi.org/10.1016/j.jenvman.2026.130515}{Shi} & Adaptive spatiotemporal graph transformer for multi-site multi-step PM$_{2.5}$ forecasting & \textit{J Environ Manage} & Deep-learning model & PM$_{2.5}$ & China / India\\
\href{https://doi.org/10.3390/bios16070380}{Tsuchiya} & Seasonal and indoor-outdoor characteristics of fluorescent aerosol particles in offices & \textit{Biosensors} & Real-time sensor campaign & Bioaerosol / coarse & Japan\\
\addlinespace[0.9mm]
\multicolumn{6}{@{}l@{}}{\cellcolor{Mist}\textsf{\textbf{\textcolor{Amber}{Exposure assessment \& modelling}}}}\\[0.6mm]
\href{https://doi.org/10.1029/2025GH001621}{Solorzano-Araque} & Multiscale exposure modelling: LOTOS-EUROS + mobility surveys + particle toxicity & \textit{GeoHealth} & CTM + mobility integration & PM$_{2.5}$ / PM$_{10}$ & Colombia\\
\addlinespace[0.9mm]
\multicolumn{6}{@{}l@{}}{\cellcolor{Mist}\textsf{\textbf{\textcolor{Coral}{Burden, policy \& mitigation}}}}\\[0.6mm]
\href{https://doi.org/10.1002/hsr2.72886}{Karamzad} & Ischaemic heart disease burden in MENA 1990-2021; PM attributable fraction 28.2\% & \textit{Health Sci Rep} & GBD secondary analysis & Ambient PM & MENA\\
\href{https://doi.org/10.1371/journal.pone.0347989}{Rees} & HEPA filtration to reduce respiratory infection in care homes & \textit{PLoS One} & RCT process evaluation & Filtered PM & UK\\
\href{https://doi.org/10.1016/j.jes.2026.01.009}{Shan} & Policy-driven emission reductions and post-pandemic rebound, high-resolution inventories & \textit{J Environ Sci (China)} & Emission inventory & PM emissions & China\\
\addlinespace[0.9mm]
\multicolumn{6}{@{}l@{}}{\cellcolor{Mist}\textsf{\textbf{\textcolor{Slate}{Other clinical endpoints}}}}\\[0.6mm]
\href{https://doi.org/10.3390/dj14070440}{Lupita} & PM$_{2.5}$, preventive-care capacity and caries/edentulism burden across the EU & \textit{Dent J} & Ecological panel & PM$_{2.5}$ & EU-27\\
\href{https://doi.org/10.1016/j.jdin.2026.06.008}{Park} & Long-term PM exposure and risk of androgenetic alopecia & \textit{JAAD Int} & Nationwide cohort & PM$_{2.5}$ / PM$_{10}$ & South Korea\\
\href{https://doi.org/10.3171/2026.3.PEDS2615}{Sun} & Environmental exposure and area deprivation vs. paediatric hydrocephalus outcomes & \textit{J Neurosurg Pediatr} & Registry cohort + geospatial & PM$_{2.5}$ (satellite-derived) & USA\\
\addlinespace[0.9mm]
\bottomrule
\end{longtable}}


\vspace{4mm}
\begin{keybox}
{\sffamily\bfseries\footnotesize\color{Deep}Provenance and method}\\[1.2mm]
{\fontsize{7.2}{8.8}\selectfont
Entry window \textbf{2026-08-26 to 2026-08-26}, a single day of deposits.
\textbf{Sources.} PubMed E-utilities on \texttt{[Date - Entry]}, run as two separate
queries (health axis 11, instrumentation axis 0). Europe PMC \texttt{CREATION\_DATE} 12.
Crossref by-ISSN across the tracked non-PubMed-indexed titles, 33 deposits after
front-matter suppression. OpenAlex remains broken and was not queried. arXiv contributed
nothing new. \textbf{Consensus} was run on the low-cost-sensor calibration axis and
returned six 2026 records, all out of window --- reported separately in the
\textit{Coverage backfill} section and excluded from every count, figure and register row
in this issue.
\textbf{Screening.} 56 raw records, 52 unique after PMID/DOI/title collapse, 2 already in
\texttt{seen.json}, \textbf{14 retained and 36 rejected} with reason logged. Rejections
were dominated by non-particle atmospheric science (PBL parameterisation, cloud
microphysics, ammonia, BVOC flux, ionospheric TEC), non-airborne chemical exposure
(mycotoxins, organophosphates, phenanthrene metabolites) and \textbf{eleven Elsevier
retraction notices deposited on a single day}.
\textbf{Two abstract-free deposits} are carried at tier C on title and journal alone and
are labelled as such in the digest.
\textbf{Effect estimates.} Nine estimates from three studies. Only the six Mao et al.\
component IRRs are mutually comparable; the COPD odds ratios use a different contrast and
the homocysteine estimate is a mediator--outcome association. See the figure caption.
\textbf{Harvester change made this run.} \textit{Meas Sci Technol} (0957-0233) and
\textit{Aerosol Sci Technol} (0278-6826) were added to the tracked-ISSN list, since both
carried 2026 low-cost-sensor calibration work that no leg of this pipeline had seen.
\textbf{Labels.} Subtopic, design, particle-metric and geography are single-label by
dominant contribution --- a judgement, not a controlled vocabulary. All DOI links were
resolved before build.\par}
\end{keybox}

\end{document}
